\section{Gram operators of connected regular regions}%
\label{sec:peps_regular_region_gram}

This section computes the Gram operator of the open tensor of a connected
region with regular bonds, and introduces coordinates on its bond labels
adapted to a spanning tree of the region.

\begin{lemma}[Expansion of the Gram matrix of a region]%
    \label{thm:peps_open_region_gram_expansion}
    \lean{TNLean.PEPS.sum_star_regionIncidentWeight_mul,
        TNLean.PEPS.openRegionCoefficientMatrix_gram_apply,
        TNLean.PEPS.regularOpenRegionMatrix_gram_apply,
        TNLean.PEPS.regularRegionIncidentConfigEquiv,
        TNLean.PEPS.regionIncidentBoundaryLabel_regular_iff,
        TNLean.PEPS.regularOpenRegionMatrix_gram_apply_group}
    \leanok
    \uses{def:peps_open_region_contraction, def:peps_regular_site_gram}
    For any finite region $R$, the entry $(\mu,\nu)$ of
    $T_R^\dagger T_R$ is the sum over pairs of incident-edge assignments
    with crossing labels $\mu,\nu$ of
    \begin{align}
        \prod_{v\in R}\sum_s
        \overline{a_v(\eta|_{E_v};s)}a_v(\theta|_{E_v};s).
    \end{align}
    This identity requires no isometry assumption.
\end{lemma}
\begin{proof}\leanok
    Expanding both copies of the region tensor, interchanging the finite
    sums, and distributing the independent physical sums over the
    vertices turns each factor into the corresponding site Gram kernel.
\end{proof}

\begin{lemma}[Synchronization of group labels in a connected region]%
    \label{thm:peps_region_group_synchronization}
    \lean{TNLean.PEPS.regionLabels_eq_of_internalEdge_eq,
        TNLean.PEPS.regionLabels_constant_iff_internalEdge_eq,
        TNLean.PEPS.IsRegionLabelCompatible,
        TNLean.PEPS.IsRegionLabelCompatible.internalEdge_eq,
        TNLean.PEPS.IsRegionLabelCompatible.exists_common_label,
        TNLean.PEPS.IsRegionLabelCompatible.exists_translation,
        TNLean.PEPS.isRegionLabelCompatible_iff_exists_translation,
        TNLean.PEPS.sum_regionLabelCompatible_eq_sum_translation}
    \leanok
    Suppose the induced graph on $R$ is nonempty and connected.
    The compatibility equations
    $\eta_e=q_v\theta_e$ at every incident leg hold precisely when
    $q_v=g$ at every vertex and $\eta=g\theta$, for one $g\in G$.
    Thus the sum over compatible vertex translations reduces to
    the sum over simultaneous translations.
\end{lemma}
\begin{proof}\leanok
    For an internal edge with endpoints $u,v$, the two equations give
    $q_u\theta_e=q_v\theta_e$, hence $q_u=q_v$.
    Equality propagates along paths. Every incident edge has an
    endpoint in $R$, so its label is translated by this common element.
    Nonemptiness makes the assignment $g\mapsto(q_v=g)$ injective.
\end{proof}

\begin{lemma}[Counting free internal regular bonds]%
    \label{thm:peps_region_internal_bond_count}
    \lean{TNLean.PEPS.sum_boundaryRestriction_eq_card,
        TNLean.PEPS.sum_translated_boundaryRestriction_eq_card,
        TNLean.PEPS.sum_boundaryRestriction_translations_eq_card,
        TNLean.PEPS.regionInternalEdgeEquiv,
        TNLean.PEPS.card_regionIncident_not_boundary,
        TNLean.PEPS.sum_regionBoundary_translations_eq_card}
    \leanok
    Fix the crossing configurations $\mu,\nu$ and a simultaneous
    translation $g$. Pairs of incident configurations satisfying
    $\eta=g\theta$ and these prescribed boundary values number
    $|G|^{I_R}$ if $\mu=g\nu$, and zero otherwise, where $I_R$ is
    the number of edges with both endpoints in $R$.
\end{lemma}
\begin{proof}\leanok
    The second configuration has one free label on each internal
    edge, and fixes the first configuration by translation.
    An incident edge that is not a crossing edge has both endpoints
    in $R$, which identifies these free indices with the internal bonds.
\end{proof}

\begin{lemma}[Gram operator of a connected regular region]%
    \label{thm:peps_connected_regular_region_gram}
    \lean{TNLean.PEPS.prod_regularSiteGram_eq_sum_compatible,
        TNLean.PEPS.sum_regularBoundary_translation_eq_card_mul_projector,
        TNLean.PEPS.regularRegionGramScalar,
        TNLean.PEPS.regularOpenRegionMatrix_gram_of_connected}
    \leanok
    \uses{thm:peps_open_region_gram_expansion, thm:peps_region_group_synchronization, thm:peps_region_internal_bond_count}
    Let $R$ induce a nonempty connected graph, with $V_R$ vertices
    and $I_R$ internal edges. Suppose the local Gram kernels are
    $A_v^\dagger A_v=c_vP_v$. For unnormalized regular bonds and
    no inserted bond operators, the open-region tensor satisfies
    \begin{align}
        T_R^\dagger T_R=\kappa_R P_{\partial R},\qquad
        \kappa_R=\left(\prod_{v\in R}c_v\right)
        \frac{|G|^{I_R+1}}{|G|^{V_R}}.
    \end{align}
    The boundary may be empty. In particular, positive local Gram
    factors give an isometry on the invariant boundary space after
    dividing by $\sqrt{\kappa_R}$, also when the region contains cycles.
    This is the untwisted connected-region contraction underlying
    \cite[Lemma~6.2 and proof of Theorem~6.9]{Schuch2010PEPS}.
\end{lemma}
\begin{proof}\leanok
    \uses{thm:peps_regular_site_gram}
    Expanding the Gram matrix into local pairings contributes
    $|G|^{-V_R}$ from their averages together with a sum over vertex group labels. Compatibility
    reduces this to one common translation. For each such translation,
    the internal labels contribute $|G|^{I_R}$. The remaining boundary
    translation sum is $|G|P_{\partial R}$.
\end{proof}

\begin{lemma}[Cycle exponent of a connected region]%
    \label{thm:peps_regular_region_cycle_exponent}
    \lean{TNLean.PEPS.card_inducedRegion_edgeSet,
        TNLean.PEPS.orderedEdgeEquivEdgeSet,
        TNLean.PEPS.inducedRegionEdgeEquiv,
        TNLean.PEPS.region_card_le_internalEdge_card_add_one,
        TNLean.PEPS.internalEdge_card_add_one_eq_card_of_isTree,
        TNLean.PEPS.regionGroupOrder_pow_div_pow_eq_cycleRank,
        TNLean.PEPS.regularRegionGramScalar_eq_cycleRank,
        TNLean.PEPS.regularRegionGramScalar_eq_prod_of_isTree}
    \leanok
    \uses{thm:peps_connected_regular_region_gram}
    In a nonempty connected region, $V_R\leq I_R+1$. Thus
    \begin{align}
        \frac{|G|^{I_R+1}}{|G|^{V_R}}
        =|G|^{I_R+1-V_R}.
    \end{align}
    The exponent is the cycle rank. For a tree it vanishes, so that
    the region Gram factor is just the product of the local factors.
\end{lemma}
\begin{proof}\leanok
    The internal edges are exactly the edges of the induced graph.
    A connected graph with $V_R$ vertices has at least $V_R-1$ edges,
    with equality for a tree. Since $|G|\ne0$, subtracting these
    nonnegative exponents gives the ratio identity.
\end{proof}

\begin{lemma}[Physical implementation of a supported canonical unitary]%
    \label{thm:peps_regular_physical_unitary_transport}
    \lean{TNLean.PEPS.exists_positive_regionPhysicalProductMatrix_gram,
        TNLean.PEPS.exists_unitary_regionPhysicalProductMatrix_mul_of_projector_commute}
    \leanok
    \uses{thm:peps_region_physical_maps}
    Let $A_v$ be the physical map of a regular $G$-isometric site,
    with arbitrary finite physical dimension, and let $P_v$ be its
    regular invariant projector. On a finite region $R$, put
    $T=\bigotimes_{v\in R}A_v$ and $P=\bigotimes_{v\in R}P_v$.
    There is $c>0$ such that $T^\dagger T=cP$, derived from local
    isometry alone. If a unitary $U$ on the canonical half-edge coordinate space
    satisfies $UP=PU$, then there is a unitary $W$ on the original
    physical region satisfying
    \begin{align}
        WT & =TU.
    \end{align}
    No surjectivity of the site maps is required. This auxiliary
    form of physical accessibility extends the operation to any unused
    physical directions; compare \cite[Lemma~6.3 and \S6.2]{Schuch2010PEPS}.
\end{lemma}
\begin{proof}\leanok
    \uses{thm:peps_regular_site_gram}
    Local isometry gives $A_v^\dagger A_v=c_vP_v$, with $c_v>0$.
    Products and adjoints of the product physical maps are computed
    site by site, so $T^\dagger T=cP$, where $c=\prod_{v\in R}c_v$.
    Hence
    \begin{align}
        (TU)^\dagger(TU) & =cU^\dagger PU=cP=T^\dagger T.
    \end{align}
    Equal Gram matrices define an isometry between the two ranges,
    which extends to a unitary on the finite-dimensional physical space.
\end{proof}

\begin{definition}[Spanning-tree coordinates for regular half-edge labels]%
    \label{def:peps_regular_region_coordinates}
    \lean{TNLean.PEPS.RootedGroupLabels,
        TNLean.PEPS.rootedGroupLabelsEquiv,
        TNLean.PEPS.rootedTreeGradient,
        TNLean.PEPS.rootedTreeGradientEquiv,
        TNLean.PEPS.RegularRegionCoordinates,
        TNLean.PEPS.regularRegionCoordinatesEquiv,
        TNLean.PEPS.regularRegionCoordinatesEquiv_symm_internal_tail,
        TNLean.PEPS.regularRegionCoordinatesEquiv_symm_internal_head,
        TNLean.PEPS.regularRegionCoordinatesEquiv_symm_boundary_tail,
        TNLean.PEPS.regularRegionCoordinatesEquiv_symm_boundary_head}
    \leanok
    \uses{def:peps_open_region_contraction}
    Let $G$ be a finite group and $R$ a finite region. Choose a spanning
    tree $T$ of its induced graph and a root $o\in R$. Write $I$ for the
    internal edges and $\partial R$ for the crossing edges, and orient
    each internal edge $e=(u,v)$ by the ordering of its endpoints.
    The tree gradient $r\mapsto(r_vr_u^{-1})_{e\in T}$ is a bijection
    from vertex labels satisfying $r_o=1$ to arbitrary labels on $T$.
    Such normalized vertex labels are themselves in bijection with
    arbitrary labels on $R\setminus\{o\}$.

    Physical labels $\alpha_{v,e}\in G$, one for each incident half-edge
    at every $v\in R$, are in bijection with coordinates
    $y\in G^{\partial R}$, $a\in G^I$, normalized vertex labels $r$,
    and $z\in G^{I\setminus T}$. The reconstruction is
    \begin{align}
        \alpha_{v,e} & =r_vy_e
                     &         & (e\in\partial R,\ v\in R),
        \label{eq:peps_region_boundary_coordinates}         \\
        \alpha_{u,e} & =r_ua_e
                     &         & (e=(u,v)\in I),\notag      \\
        \alpha_{v,e} & =
        \begin{cases}
            r_va_e,    & e\in T,            \\
            r_vz_ea_e, & e\in I\setminus T.
        \end{cases}
        \label{eq:peps_region_internal_coordinates}
    \end{align}
    This is a reversible change of regular basis labels associated with
    the accessible virtual systems and blocking operations in
    \cite[\S6.2--6.3]{Schuch2010PEPS}.
\end{definition}
\begin{proof}\leanok
    \uses{thm:peps_regular_region_cycle_exponent}
    If two normalized vertex labellings have the same tree gradient,
    their relative label $r_v^{-1}s_v$ is constant along every tree edge.
    Connectedness and the root normalization give $r=s$. A tree has
    $|R|-1$ edges, so the domain and codomain of the gradient each have
    $|G|^{|R|-1}$ elements. Its injectivity therefore gives bijectivity.

    From $\alpha$, recovering $r$ via the tree ratios
    $\alpha_{v,e}\alpha_{u,e}^{-1}$ and setting
    $y_e=r_v^{-1}\alpha_{v,e}$ on the boundary,
    $a_e=r_u^{-1}\alpha_{u,e}$ on each internal edge, and
    $z_e=r_v^{-1}\alpha_{v,e}a_e^{-1}$ outside the tree
    recovers every original label upon substitution into
    (\ref{eq:peps_region_boundary_coordinates})--(\ref{eq:peps_region_internal_coordinates}).
    Conversely, the reconstructed tree
    ratios equal $r_vr_u^{-1}$, so uniqueness recovers $r$, and then
    the same formulas recover $y$, $a$, and $z$.
\end{proof}

\begin{lemma}[Canonical regular block in spanning-tree coordinates]%
    \label{thm:peps_regular_region_projector_coordinates}
    \lean{TNLean.PEPS.regularRegionCoordinates_rootLabels_of_tree_compatibility,
        TNLean.PEPS.regularProjectorOpenRegionMatrix_coordinates,
        TNLean.PEPS.regularProjectorOpenRegionMatrix_exists_normalized_ancilla}
    \leanok
    \uses{def:peps_regular_region_coordinates}
    Let $R$ have a spanning tree $T$ rooted at $o$, and let every local
    tensor be the simultaneous regular averaging projector. Write
    $\alpha(y,a,r,z)$ for the reconstructed half-edge labels in
    (\ref{eq:peps_region_boundary_coordinates})--(\ref{eq:peps_region_internal_coordinates}).
    The open-region coefficient matrix $M_R$ satisfies
    \begin{align}
        M_R\bigl(\alpha(y,a,r,z),\theta\bigr)
         & = |G|^{1-|R|}P_{\partial R}(y,\theta)
        \mathbf 1_{\{z=1\}}.
        \label{eq:peps_region_projector_coordinates}
    \end{align}
    No nonempty-boundary assumption is required.

    Let $N$ be the number of pairs $(a,r)$, with $r_o=1$, and put
    $\phi(a,r,z)=N^{-1/2}\mathbf 1_{\{z=1\}}$. In the coefficient form of the regular blocking construction
    of \cite[\S6.2--6.3]{Schuch2010PEPS}, $\|\phi\|^2=1$,
    and (\ref{eq:peps_region_projector_coordinates}) factors as
    $\gamma P_{\partial R}(y,\theta)\phi(a,r,z)$, where
    $\gamma=|G|^{1-|R|}\sqrt N>0$.
\end{lemma}
\begin{proof}\leanok
    \uses{thm:peps_regular_open_region_projector}
    Expanding the original internal-bond contraction into a sum over bond
    labels $\eta$ and vertex translations $q$ gives
    $r_ua_e=q_u\eta_e$ and $r_va_e=q_v\eta_e$ on each tree edge
    $e=(u,v)$ by compatibility.
    Thus $q_vq_u^{-1}=r_vr_u^{-1}$. The uniqueness of rooted tree
    gradients gives $q_v=r_vx$, where $x=q_o$; this conclusion already
    follows from compatibility on the tree alone.

    Compatibility on the remaining internal edges forces $z_e=1$.
    The bond labels are then uniquely determined by $x$, and boundary
    compatibility is precisely $y=x\theta$. Consequently the coefficient
    is $|G|^{-|R|}\sum_x\mathbf 1_{\{y=x\theta,\ z=1\}}$,
    which is (\ref{eq:peps_region_projector_coordinates}). There are
    exactly $N$ nonzero ancillary coefficients, each of modulus
    $N^{-1/2}$, proving the stated normalization.
\end{proof}

\begin{lemma}[Isometry of a connected regular block]%
    \label{thm:peps_connected_regular_region_isometry}
    \lean{TNLean.PEPS.regularOpenRegionMatrix_invariant,
        TNLean.PEPS.isGIsometric_regularOpenRegionMatrix_of_connected}
    \leanok
    \uses{thm:peps_connected_regular_region_gram}
    Let every site be regular $G$-isometric, and let $R$ induce a nonempty
    connected graph. The open-region map $T_R$ is $G$-isometric for
    simultaneous regular translation on the crossing bonds. Its derived
    positive factor is
    \begin{align}
        c_R=\left(\prod_{v\in R}c_v\right)
        |G|^{I_R+1-|R|},
    \end{align}
    where $I_R$ is the number of internal edges. This includes blocks with
    internal cycles and blocks with no crossing bonds. It is the
    finite-region form of the stability and blocking argument in
    \cite[Lemma~6.2 and \S6.2]{Schuch2010PEPS}, with the positive-factor
    convention recorded in \cite{gap:rmp_peps_quantum_double_g_isometry}.
\end{lemma}
\begin{proof}\leanok
    \uses{thm:peps_regular_open_region_projector}
    Translating all virtual labels in the contraction proves boundary
    invariance. The proved connected-region Gram formula gives
    $T_R^\dagger T_R=c_R\Pi_{\partial R}$ with $c_R>0$.
    Dividing the adjoint by $c_R$ gives a left inverse on the invariant
    boundary space. The same identity gives the required inner-product
    identity there.
\end{proof}

\begin{lemma}[Injectivity of a connected regular block]%
    \label{thm:peps_connected_regular_region_injectivity}
    \lean{TNLean.PEPS.isGInjective_regularOpenRegionMatrix_of_connected}
    \leanok
    \uses{def:peps_g_injective}
    Let $R$ be a nonempty connected finite region whose site tensors
    are $G$-injective for the simultaneous regular action on their
    incident legs. Then the open tensor of $R$ is $G$-injective
    for the simultaneous regular action on its crossing legs.
    Internal cycles and an empty boundary are allowed.
    This is the regular-representation connected-region consequence of
    \cite[Lemma~5.2 and the following observation]{Schuch2010PEPS};
    the argument does not assume local isometry.
\end{lemma}
\begin{proof}\leanok
    \uses{thm:peps_regular_region_projector_coordinates, thm:peps_regular_open_region_accessibility}
    Applying the product of the local inverses to the open tensor yields
    its canonical averaging-projector contraction.
    Choosing a spanning tree and a root, in the coordinates of
    Lemma~\ref{thm:peps_regular_region_projector_coordinates} the row
    with identity reference, normalized vertex, and residual cycle labels
    is $|G|^{1-|R|}\Pi_{\partial R}$. This scalar is nonzero, and
    $\Pi_{\partial R}$ fixes every invariant boundary vector. Thus a
    vanishing physical output implies a vanishing invariant input.
    Boundary invariance follows by translating every internal label
    in the contraction.
\end{proof}

\begin{lemma}[Removing bond operators on a spanning tree]%
    \label{thm:peps_regular_region_tree_gauge}
    \lean{TNLean.PEPS.regularRegionTreeGauge_root,
        TNLean.PEPS.regularRegionTreeGauge,
        TNLean.PEPS.regularRegionTreeGauge_gradient,
        TNLean.PEPS.regularRegionGaugeResidual,
        TNLean.PEPS.regularRegionGaugeResidual_transport,
        TNLean.PEPS.regularRegionGaugeResidual_eq_one_of_tree,
        TNLean.PEPS.regularRegionTreeCycleResidual,
        TNLean.PEPS.regularRegionGaugeInternalLabels,
        TNLean.PEPS.regularRegionGaugeInternalLabels_apply,
        TNLean.PEPS.regularRegionGaugePhysicalLabels,
        TNLean.PEPS.regularRegionGaugePhysicalLabels_apply,
        TNLean.PEPS.regularRegionBoundaryTransport,
        TNLean.PEPS.regularRegionBoundaryTransport_tail,
        TNLean.PEPS.regularRegionBoundaryTransport_head,
        TNLean.PEPS.regularTwistedLabels_gauge}
    \leanok
    \uses{def:peps_regular_region_coordinates,
        def:peps_regular_twisted_region}
    Let $G$ be finite, choose a spanning tree $T$ of the region $R$,
    and fix a root $o$. Assign a group operator $u_e$ to each oriented
    edge $e=(s,t)$. There is a unique vertex gauge $k$ with $k_o=1$
    and $k_tk_s^{-1}=u_e$ on $T$. The internal residual operators
    $\omega_e=k_t^{-1}u_ek_s$ therefore equal the identity on $T$;
    their values on the remaining internal edges are unrestricted.

    Internal references and physical half-edge labels transform by the invertible
    shifts $\eta'_e=k_s^{-1}\eta_e$ and
    $\alpha'_{v,e}=k_v^{-1}\alpha_{v,e}$. Crossing labels transform by the
    invertible map
    \begin{align}
        \theta'_e & =
        \begin{cases}
            k_s^{-1}\theta_e,    & s\in R,\ t\notin R, \\
            k_t^{-1}u_e\theta_e, & s\notin R,\ t\in R.
        \end{cases}
        \label{eq:peps_tree_boundary_transport}
    \end{align}
    On every internal edge,
    $\omega_e\eta'_e=k_t^{-1}u_e\eta_e$.
    More generally, for any ambient vertex labels $K$, set
    $\widetilde u_e=K_t^{-1}u_eK_s$ and
    $\widetilde\eta_{v,e}=K_s^{-1}\eta_{v,e}$.
    If $D_u$ inserts $u_e$ at the head of each edge and leaves its
    tail label unchanged, then
    $D_{\widetilde u}(\widetilde\eta)_{v,e}
        =K_v^{-1}D_u(\eta)_{v,e}$.
    These are group-coordinate identities associated with the reversible
    regular blocking operations in \cite[\S6.2--6.3]{Schuch2010PEPS}.
\end{lemma}
\begin{proof}\leanok
    \uses{def:peps_regular_region_coordinates,
        def:peps_regular_twisted_region}
    Inverting the tree-gradient bijection at $(u_e)_{e\in T}$ yields
    the normalized $k$, so on a tree edge
    $k_t^{-1}u_ek_s=k_t^{-1}(k_tk_s^{-1})k_s=1$.
    The reference identity follows from
    $(k_t^{-1}u_ek_s)(k_s^{-1}\eta_e)=k_t^{-1}u_e\eta_e$.
    Each map is multiplication by a fixed group element at each
    coordinate, hence is reversible. At a head endpoint the native
    twisted-label identity is
    $(K_t^{-1}u_eK_s)(K_s^{-1}\eta_{t,e})
        =K_t^{-1}u_e\eta_{t,e}$;
    at a tail endpoint both expressions are $K_s^{-1}\eta_{s,e}$.
\end{proof}
